\documentclass[11pt]{article}
\usepackage[margin=0.82in]{geometry}
\usepackage{amsmath,amssymb}
\usepackage{graphicx}
\usepackage[colorlinks=true,linkcolor=blue,urlcolor=blue,citecolor=blue]{hyperref}
\setlength{\parindent}{0pt}
\setlength{\parskip}{0.56em}

\newcommand{\R}{\mathbb R}
\newcommand{\I}{\mathbb I}

\title{Two optimal Lorentz endpoint questions\\
\large Affirmative answers in arXiv:2010.05297}
\author{Run \texttt{fa\_banach\_001} --- \texttt{agent\_lane\_03}}
\date{12 August 2026}

\begin{document}
\maketitle

\textbf{Status.} \texttt{literature\_already\_answered}.  This packet
identifies and checks a later affirmative answer to both questions.  It does
not claim a new theorem.

\section*{The two source questions}

Spector and Van Schaftingen, \emph{Optimal embeddings into Lorentz spaces
for some vector differential operators via Gagliardo's lemma},
arXiv:1811.02691, ask whether every homogeneous constant-coefficient
elliptic canceling operator has the sharp Lorentz endpoint.  Their Questions
1.1 and 1.2 ask, respectively, whether
\begin{align}
 \|u\|_{L^{n/(n-1),1}(\R^n)}
   &\le C\|A(D)u\|_{L^1(\R^n)} &&(\operatorname{ord}A=1), \tag{Q1}\\
 \|D^{k-1}u\|_{L^{n/(n-1),1}(\R^n)}
   &\le C\|A(D)u\|_{L^1(\R^n)} &&(\operatorname{ord}A=k\ge2). \tag{Q2}
\end{align}
The dimensions in (Q1) are $n\ge3$ and in (Q2) $n\ge2$.

\begin{center}
\includegraphics[width=0.97\textwidth,trim=0 330 0 0,clip]{figures/source_page4.png}
\end{center}

\section*{Exact later answer}

Stolyarov, \emph{Hardy--Littlewood--Sobolev inequality for $p=1$},
arXiv:2010.05297v3 (2021), proves the required estimate in all orders.  His
abstract records as a particular case that, for every canceling elliptic
differential operator $A$ of order $m$,
\begin{equation}
 \boxed{\quad
 \|\nabla^{m-1}u\|_{L^{d/(d-1),1}(\R^d)}
 \lesssim \|Au\|_{L^1(\R^d)}.
 \quad} \tag{S}
\end{equation}
Taking $m=1$ gives (Q1), since $\nabla^0u=u$.  Taking $m=k\ge2$
gives (Q2), with equivalent finite-dimensional norms on the full
$(k-1)$st derivative tensor.  Stolyarov's statement holds for every $d\ge2$;
it therefore covers all dimensions asked in the source.

\section*{The theorem-to-question bridge}

Here is the mechanism inside the later paper, included to rule out a merely
notational match.  For an order-$m$ elliptic operator $A:V\to E$, Stolyarov's
Example 1.4 assigns
\[
 \Omega(\zeta)=\operatorname{Im}A(\zeta),\qquad \zeta\in S^{d-1}.
\]
The source cancellation hypothesis is exactly
$\bigcap_{\zeta\in S^{d-1}}\Omega(\zeta)=\{0\}$.  Hence Corollary 1.7 of
the later paper, with $\alpha=1$, gives
\[
 \|\I_1(Au)\|_{L^{d/(d-1),1}}
 \lesssim \|Au\|_{L^1}.
\]
Ellipticity supplies a smooth homogeneous left inverse of the symbol:
for $\xi\ne0$, one may use
\[
 A(\xi)^\dagger=(A(\xi)^*A(\xi))^{-1}A(\xi)^*.
\]
Consequently each component of $D^{m-1}u$ is a degree-zero smooth Fourier
multiplier applied to $\I_1(Au)$.  Such Calder\'on--Zygmund multipliers are
bounded on $L^{d/(d-1),1}$, whose first exponent is strictly between $1$
and $\infty$.  This proves (S), exactly with the source's domain, operator
hypotheses, derivative order, and Lorentz index.

\section*{Proof intuition}

An unconstrained $L^1$ datum can concentrate into a point mass, which is the
endpoint obstruction for the Riesz potential.  Cancellation forbids that
common point-mass direction in the symbol images.  Stolyarov proves that this
is precisely enough to recover the strongest Lorentz target for $\I_1$.
Ellipticity then reconstructs one fewer derivative of $u$ from that potential
using only order-zero singular integrals.  Thus cancellation removes the
$L^1$ concentration obstruction, while ellipticity turns the potential
estimate into the desired differential estimate.

\section*{Scope and provenance}

This is a complete affirmative resolution of both displayed source
questions.  The supporting paper was posted on 11 October 2020 and the
bundled PDF is arXiv:2010.05297v3, dated 5 February 2021.  Its bibliography
includes the source paper, and its abstract explicitly states the full
higher-order estimate.  The bundled one-page source excerpt contains both
questions; the four-page supporting excerpt contains the abstract, Example
1.4, the cancellation definition, and Corollary 1.7.

{\small
\textbf{References.}
[1] D. Spector and J. Van Schaftingen, \emph{Optimal embeddings into Lorentz
spaces for some vector differential operators via Gagliardo's lemma},
arXiv:1811.02691, Questions 1.1--1.2.
[2] D. Stolyarov, \emph{Hardy--Littlewood--Sobolev inequality for $p=1$},
arXiv:2010.05297v3; Sb. Math. 213 (2022), 844--889.
}

\end{document}
